\newpage
\thispagestyle{empty}
\begin{center}
  \textbf{\MakeUppercase{Resumo}}
\end{center}

\noindent
Esta dissertação investiga como problemas de otimização podem ser compreendidos
dentro de um arcabouço composicional e diagramático unificado, fundamentado na
teoria das categorias. A motivação central é a observação de que estruturas
matemáticas essenciais em otimização --- relações, funções convexas, formas
quadráticas, poliedros e programas lineares --- podem ser interpretadas como
morfismos em categorias nas quais a composição corresponde à eliminação ou
minimização de variáveis internas. A partir desta perspectiva, a otimização
emerge não como uma coleção de técnicas isoladas, mas como uma operação
algébrica inerente à estrutura relacional.
O desenvolvimento segue uma progressão incremental, partindo das relações
binárias clássicas enriquecidas por quantais de Lawvere. A cada etapa é
introduzido um único gerador novo e um conjunto finito de axiomas. A Álgebra
Quadrática Gráfica (GQA) representa problemas de minimização da norma
$\ell_2$ sobre subespaços afins; a Álgebra Poliedral enriquece o arcabouço com
restrições de desigualdade linear; e a Álgebra Gráfica para Programação Linear
Paramétrica (GLP) constitui a contribuição central da tese, na qual um programa
linear é tratado como um morfismo compositável com outros blocos de otimização.
Para cada álgebra são demonstrados resultados de solidez e completude com
respeito às categorias semânticas correspondentes. Um resultado central mostra
que toda função convexa afim por partes admite representação como função valor
de um programa linear, e que essa equivalência se manifesta diagramaticamente
como uma forma normal canônica, eliminando no nível sintático a distinção entre
um LP e uma função convexa. O cálculo diagramático para GLP desempenha, para a
programação linear, o mesmo papel que a Álgebra Linear Gráfica desempenha para
matrizes: fornece uma linguagem de formas normais, regras de reescrita e
teoremas de completude dentro dos quais argumentos de otimização podem ser
conduzidos simbolicamente, antes e independentemente de qualquer computação
numérica.

\vspace{1em}
\noindent
Palavras-chave: teoria das categorias; diagramas de corda; programação linear;
otimização convexa; categorias monoidais simétricas; relações ponderadas;
quantais; álgebra gráfica; formas normais; completude.
\clearpage
